\documentclass[10pt,a4paper]{amsart}
\usepackage[margin=19mm]{geometry}
\usepackage{amsmath,amssymb,mathtools}
\usepackage[scr=rsfs,cal=euler]{mathalfa}
\usepackage{xfrac}
\usepackage[unicode,hidelinks,bookmarks=false]{hyperref}
\newcommand{\ie}{i.e.\ }
\newcommand{\cf}{cf.\ }
\newcommand{\resp}{respectively\ }
\newcommand{\Subsection}{\subsection}
\providecommand{\nobreakdash}{\nobreak\hbox{-}}
\setlength{\emergencystretch}{2em}
\multlinegap0pt
\usepackage{bm}
\usepackage{bbm}
\usepackage{dsfont}
\usepackage{xspace}
\usepackage{cancel}
\usepackage{mathtools}
\usepackage[shortlabels]{enumitem}
\usepackage{amsthm}
\makeatletter
\@ifundefined{thm}{\newtheorem{thm}{Thm}}{}
\@ifundefined{assumption}{\newtheorem{assumption}[thm]{Assumption}}{}
\@ifundefined{corollary}{\newtheorem{corollary}[thm]{Corollary}}{}
\makeatother
\newcommand\bN{\mathbb{N}}
\newcommand\bR{\mathbb{R}}
\newcommand\cA{\mathcal{A}}
\newcommand\cF{\mathcal{F}}
\title[Well-Posedness for Cauchy Problems with Singular Time-Measur]{Well-Posedness for Cauchy Problems with Singular Time-Measurable Pseudo-Differential Operators in Quasi-decreasing Weighted $\mathrm{L}\_2$-Spaces}
\author{Jae-Hwan Choi, Ildoo Kim}
\date{Source: arXiv:2608.19717v1 (CC BY 4.0)}
\begin{document}
\maketitle

\noindent\emph{Source and licence.} Well-Posedness for Cauchy Problems with Singular Time-Measurable Pseudo-Differential Operators in Quasi-decreasing Weighted $\mathrm{L}_2$-Spaces. Version arXiv:2608.19717v1 (CC BY 4.0), distributed under the Creative Commons Attribution 4.0 International licence, as recorded in the arXiv licence metadata for this exact version and shown on the abstract page https://arxiv.org/abs/2608.19717v1. Full rights notice: licenses/arxiv-2608.19717-CC-BY-4.0.txt.\par\smallskip

\noindent\textit{Editorial scope.} Counted unit: the Corollary whose source label is \texttt{dense cor} in the article (source0.tex lines 1919--1932, source label \texttt{dense cor}), delivered together with the article's own complete original proof of it (source0.tex lines 1933--2017, one \texttt{proof} environment of 85 lines). No mathematics is written, repaired or re-derived: statement and proof are the article's own text line for line. Required context retained uncounted: weight as (lines 694--707); 20260811 01 (lines 1384--1392); 20260401 01 (lines 1748--1753). Every definition, display and label of those lines is retained unchanged. Omitted from this package and not redistributed: 1--693, 708--1383, 1393--1747, 1754--1918, 2018--2847. Nothing inside a retained excerpt is omitted or altered: every retained body is delivered exactly as the article's lines are written, so no omission receipt of any kind accompanies this package. Citations: the retained excerpts carry no citation command at all, so no bibliography entry of the article is transcribed and no reference is redistributed. Reference adaptation: none was needed and none was made; every reference inside the delivered unit resolves inside it, so no unresolved reference or citation is produced. Printed slips kept literal: no printed word, symbol, hypothesis or number of the counted statement or of its proof was altered. Layout and class: the article's own class is \texttt{amsart} with packages including enumerate, amsmath, amssymb, amsthm, mathrsfs, esint, mathtools, hyperref, bm, xcolor; the wrapper is the lane's pinned \texttt{amsart} editorial wrapper at 10pt on A4, which restates the theorem-like environments the retained text needs and adds the article's own macro definitions that the retained text needs (\texttt{\textbackslash bN}, \texttt{\textbackslash bR}, \texttt{\textbackslash cA}, \texttt{\textbackslash cF}), with extra wrapper packages (bm, bbm, dsfont, xspace, cancel, mathtools, enumitem). The delivered numbering is the unit's own. Assets: no retained span contains a figure environment, a graphics-inclusion command, a PDF-page inclusion, a table or a TikZ picture; no image file is copied into this package, the package declares no asset dependency and no image file is redistributed with it. Licence: version arXiv:2608.19717v1 of this article is distributed under the Creative Commons Attribution 4.0 International licence, as recorded in the arXiv licence metadata for this exact version and shown on the abstract page https://arxiv.org/abs/2608.19717v1. A full rights notice is delivered at licenses/arxiv-2608.19717-CC-BY-4.0.txt. Characters: every retained excerpt is pure ASCII, so the delivered engine, XeLaTeX with its default Latin Modern text font, reports no missing character.
\begin{assumption}[Quasi-decreasing Weight]
										\label{weight as}
Let $w$ be a real-valued measurable function defined a.e. on the interval $(0,T)$. We assume that $w$ satisfies the following conditions:
\begin{enumerate}[(i)]
\item The weight is strictly positive almost everywhere:
\begin{align*}
w(t) >0 \quad \text{a.e.}~ t \in (0,T)
\end{align*}
\item There exists a constant $N_w > 0$ such that the weight is quasi-decreasing; that is,
\begin{align*}
\frac{w(t)}{w(s)} \leq N_w \quad \text{a.e.}~ 0<s<t<T.
\end{align*}
\end{enumerate}
\end{assumption}

\begin{equation}
										\label{20260811 01}
\begin{aligned}
&\int_0^t \int_{B_R} |\psi(s,\xi)|\mathrm{d}\xi \mathrm{d}s < \infty
\quad \text{for all } t \in (0,T) \text{ and } R > 0,\\
&\psi(s,\cdot)\in \mathrm{L}_{2,loc}(\bR^d)
\quad \text{for a.e. }s\in(0,T).
\end{aligned}
\end{equation}

\begin{align}
									\label{20260401 01}
\int_0^t  \left(\int_{\bR^d}  \left( |\psi(\rho,\xi)| \int_0^\rho |\cF[f(s,\cdot)](\xi) \mathrm{d}s\right)^2  \mathrm{d}\xi \right)^{1/2} \mathrm{d}\rho
+\int_0^t  \int_{\bR^d}   |\psi(\rho,\xi)| \int_0^\rho |\cF[f(s,\cdot)](\xi) \mathrm{d}s  \mathrm{d}\xi \mathrm{d}\rho
<\infty.
\end{align}

\begin{corollary}
									\label{dense cor}
Let $w$ denote a weight function satisfying Assumption \ref{weight as}, and let $\phi$ be a complex-valued function defined almost everywhere on $\bR^d$. 
Under these conditions, the scaled set
$$
\frac{1}{\sqrt{w}} \cA := \left\{ \frac{1}{\sqrt{w(t)}} \sum_{i=1}^n 1_{(a_i,b_i)}(t) f_i(x) : n \in \bN, \; 0 < a_i < b_i < T, \; f_i \in \cF^{-1}\mathrm{C}_c^\infty(\bR^d) \right\}
$$
is dense in the weighted space $\mathrm{L}_{2,loc}\left( (0,T) \times \bR^d, w(t)\mathrm{d}t\mathrm{d}x \right)$.
Then, as a direct result, assuming \eqref{20260811 01} holds, those $f$ belonging to
 $$
\mathrm{L}_{2,loc}\left( (0,T) \times \bR^d, w(t)\mathrm{d}t \mathrm{d}x \right)\cap \cF^{-1}\mathrm{L}_{1,loc}\left( (0,T) \times \bR^d \right), 
$$ 
that satisfy \eqref{20260401 01} form a dense subspace of $\mathrm{L}_{2,loc}\left( (0,T) \times \bR^d, w(t)\mathrm{d}t\mathrm{d}x \right)$.
\end{corollary}

\begin{proof}
First, observe that if a function $f$ belongs to the weighted space $\mathrm{L}_{2,loc}\left( (0,T) \times \bR^d, w(t)\mathrm{d}t\mathrm{d}x \right)$, then the product $f \sqrt{w}$ naturally resides in the standard, unweighted space $\mathrm{L}_{2,loc}\left( (0,T) \times \bR^d \right)$.
By the preceding theorem, the class $\cA$ is dense in $\mathrm{L}_{2,loc}\left( (0,T) \times \bR^d \right)$. 
Therefore, the function $f \sqrt{w}$ can be approximated by elements from $\cA$. Dividing this approximation by $\sqrt{w}$ implies that the set $\frac{1}{\sqrt{w}} \cA$ is dense in $\mathrm{L}_{2,loc}\left( (0,T) \times \bR^d, w(t)\mathrm{d}t\mathrm{d}x \right)$.

It remains to prove the second assertion. 
Let $\mathscr{D}$ denote the class of functions appearing therein. 
Since $\cA/\sqrt w$ is dense in the weighted space, it suffices to prove that
$$
\frac{1}{\sqrt w}\cA\subset\mathscr D.
$$
Fix
$$
g(s,x)
=
\frac{1}{\sqrt{w(s)}}
\sum_{i=1}^n1_{(a_i,b_i)}(s)f_i(x)
\in\frac{1}{\sqrt w}\cA,
$$
and set
$$
a_*:=\min_{1\leq i\leq n}a_i,
\qquad
b_*:=\max_{1\leq i\leq n}b_i,
\qquad
K:=\bigcup_{i=1}^n\operatorname{supp}\cF[f_i].
$$
By Assumption~\ref{weight as}, $w^{-1/2}$ is essentially bounded on
$(a_*,b_*)$. Since $\cF[f_i]\in C_c^\infty(\bR^d)$, it follows immediately
that
$$
g\in
\mathrm{L}_{2,loc}\left((0,T)\times\bR^d,w(t)\,\mathrm dt\,\mathrm dx\right)
\cap
\cF^{-1}\mathrm{L}_{1,loc}\left((0,T)\times\bR^d\right).
$$

Define
$$
G(\rho,\xi)
:=
\int_0^\rho
\left|\cF[g(s,\cdot)](\xi)\right|\,\mathrm ds.
$$
For some constant $C_g>0$,
$$
G(\rho,\xi)
\leq
C_g1_{[a_*,T)}(\rho)1_K(\xi).
$$
Thus, when $t\leq a_*$, both terms in \eqref{20260401 01} vanish. For
$t>a_*$, \eqref{20260811 01}, the local boundedness of $w$, and the
Cauchy--Schwarz inequality yield
\begin{align*}
\int_{a_*}^t
\left(
    \int_K|\psi(\rho,\xi)|^2\,\mathrm d\xi
\right)^{1/2}
\mathrm d\rho
&\leq
\left(
    \int_{a_*}^t\int_K
    |\psi(\rho,\xi)|^2\frac{\rho^2}{w(\rho)}
    \,\mathrm d\xi\,\mathrm d\rho
\right)^{1/2}
\left(
    \int_{a_*}^t\frac{w(\rho)}{\rho^2}\,\mathrm d\rho
\right)^{1/2}
<\infty.
\end{align*}
Moreover, since $K$ has finite measure,
$$
\int_{a_*}^t\int_K|\psi(\rho,\xi)|\,\mathrm d\xi\,\mathrm d\rho
\leq
|K|^{1/2}
\int_{a_*}^t
\left(
    \int_K|\psi(\rho,\xi)|^2\,\mathrm d\xi
\right)^{1/2}
\mathrm d\rho
<\infty.
$$
Together with the bound on $G$, these estimates imply \eqref{20260401 01}. 
Hence $\cA/\sqrt w\subset\mathscr D$, and the second assertion follows from the density of $\cA/\sqrt w$.
\end{proof}

\end{document}
